\documentclass{article}
\usepackage{/globals/de}
\usepackage{/globals/anki}
\ankiprefix{verwaltungskompetenzen}
\ankitopic{Verwaltungskompetenzen}
\ankishow
\begin{document}
\section{Verwaltungskompetenzen} 
Die Verwaltungskompetenzen heißt es für einen Verband (Bund oder Land) das Recht und die Pflicht zur staatlichen Gewalt für Gesetze bestimmter Materien zu haben.
\begin{description}
 \item[Gesetzesakzessorische Verwaltung] Die Exekutive führt Gesetze aus. Sowohl Landes- als auch Bundesgesetze werden \idr von den Ländern ausgeführt, weshalb es einfach zu \emph{Inkongruenz} kommt; der Verband der ein Gesetze erlassen hat ist nicht der Verband der aus ausführt.
\ankinote{gesetzesakzessorische}{Gesetzesakzessorische}{Gesetze von Exekutive ausgeführt, idr Länder}
\ankinote{inkongruenz}{Inkongruenz}{Erlasser ist nicht Ausführer}
 \item[Nicht-gesetzesakzessorische Verwaltung] Es kommt zu Verwaltungstätigkeiten, ohne dass dabei ein bestimmtes Fachgesetz ausgeführt wird. Auch diese werden \idr von den Ländern ausgeführt.
\ankinote{nicht-gesetzesakzessorische}{Nicht-gesetzesakzessorische}{Verwaltung ohne Gesetz}
\end{description}

\subsection{Ausführung von Bundesgesetzen durch die Länder}
Die Länder können Bundesgesetze auf zwei verschiedene Art und Weisen ausführen
\begin{description}
 \item[Eigene Angelegenheit] In der Regel werden Bundesgesetze als \emph{eigene Angelegenheit} ausgeführt.

Die Länder können selbst über die Verwaltungsverfahren bestimmen.

Durch Vorgaben in Bundesgesetzen kann der Bund eingreifen, von diesen kann aber abgewichen werden. Der Bund kann abweichungsfeste Vorgaben erlassen, mit Zustimmung vom Bundesrat.

Die Bundesregierung hat die Rechtsaufsicht; sie überprüfen die Rechtsmäßigkeit. Dazu können sie Beauftragte in die Ministerien entsenden.
\ankinote{eigene-angelegenheit}{Eigene Angelegenheit}{Regelfall, Länder selbst Verfahren, Bund nur wenig eingreifen. Beauftragte in Ministerien}
 \item[Im Auftrage des Bundes] Bestimmte Materien können \emph{im Auftrage des Bundes} ausgeführt werden.

Der Bund kann Verfahrensvorschriften erlassen, mit dem Bundesrat zusammen die Verwaltungsvorschriften bestimmen.

Die Bundesregierung kann Beauftrage in Behörden entsenden, Berichte und Akten verlangen.
\ankinote{im-auftrage}{Im Auftrage des Bundes}{Bestimmte Materien, Bund Verfahrensvorschriften, Beauftrage in Behörden, Berichte, Akten}
\end{description}
In beiden Fällen können die Länder selbst über die Einrichtung der Behörden bestimmen.
\ankinote{diff}{Eigene Angelegenheit vs Im Auftrage}{Grad der Rechtsaufsicht, des Eingreifens}

\subsection{Bundeseigene Verwaltung}
Der Bund führt Gesetze zu bestimmten Materien selbst aus. Diese kann unmittelbar oder mittelbar sein.\begin{description}
 \item[Bundesunmittelbare Verwaltung] Durch einen eigenen Verwaltungsunterbau.
\ankinote{unmittelbar}{Bundesmittelbare Verwaltung}{Eigenen Verwaltungsunterbau}
 \item[Mittelbare Bundesverwaltung] Durch Körperschaften oder Anstalten des öffentlichen Rechts.
\ankinote{mittelbar}{Mittelbare Bundesverwaltung}{Körperschaften oder Anstalten ÖR}
\end{description}
Die Bundesregierung erlässt allgemeine Verwaltungsvorschriften und bestimmt die Einrichtung der Behörden, ohne Bundesrat.
\ankinote{bundeseigene-verwaltung}{Bundeseigene Verwaltung}{BR Verwaltungsvorschriften, Einrichtung der Behörden, ohne BR}
\end{document}